% =============================================================================
% Paper Tables - Small LLM + UMLS KG for Differential Diagnosis
% =============================================================================

% -----------------------------------------------------------------------------
% Table 1: Main Results - Comparison with Baselines
% -----------------------------------------------------------------------------
\begin{table}[t]
\centering
\caption{Comparison with baseline methods on DDXPlus dataset (severity=2, n=1,000).
Our method (Small LLM + KG) outperforms AARLC by +10.6\% in GTPA@1 while requiring 6.6 fewer questions on average.}
\label{tab:main_results}
\begin{tabular}{lcccc}
\toprule
\textbf{Method} & \textbf{GTPA@1} & \textbf{DDR} & \textbf{DDF1} & \textbf{IL} \\
\midrule
AARLC \cite{ddxplus} & 75.4\% & 97.7\% & 78.2\% & 25.8 \\
\midrule
\multicolumn{5}{l}{\textit{Ours (Small LLM + KG)}} \\
\quad Ministral-3B & 86.4\% & 52.8\% & 47.3\% & 19.1 \\
\quad Qwen3-4B & 86.0\% & 52.7\% & 47.3\% & 19.1 \\
\quad gemma-3-4b & 86.2\% & 52.8\% & 47.3\% & 18.6 \\
\quad medgemma-1.5-4b & 86.3\% & 52.9\% & 47.5\% & 18.9 \\
\quad Phi-4-mini & 86.3\% & 52.7\% & 47.3\% & 18.6 \\
\quad Qwen3-8B & 86.2\% & 52.8\% & 47.4\% & 19.1 \\
\quad Ministral-8B & 86.2\% & 52.8\% & 47.3\% & 18.6 \\
\quad Llama-3.1-8B & 85.6\% & 52.9\% & 47.4\% & 19.2 \\
\midrule
\textbf{Average} & \textbf{86.2\%} & 52.8\% & 47.4\% & \textbf{18.9} \\
\bottomrule
\end{tabular}
\end{table}

% -----------------------------------------------------------------------------
% Table 2: Top-K Ablation Study (Joint: Symptom & Diagnosis)
% -----------------------------------------------------------------------------
\begin{table}[t]
\centering
\caption{Ablation study on the number of candidates $K$ presented to the LLM for both symptom selection and diagnosis selection. $K=3$ achieves optimal balance between accuracy and efficiency.}
\label{tab:topk_joint}
\begin{tabular}{ccc}
\toprule
\textbf{K} & \textbf{GTPA@1} & \textbf{IL} \\
\midrule
1 & 85.8\% & 19.1 \\
2 & 85.6\% & 19.4 \\
\textbf{3} & \textbf{85.6\%} & \textbf{19.3} \\
4 & 85.4\% & 19.6 \\
5 & 85.2\% & 20.8 \\
6 & 85.0\% & 21.3 \\
8 & 85.0\% & 21.7 \\
10 & 85.2\% & 21.8 \\
\bottomrule
\end{tabular}
\end{table}

% -----------------------------------------------------------------------------
% Table 3: Diagnosis Top-K Ablation (Symptom K=3 Fixed)
% -----------------------------------------------------------------------------
\begin{table}[t]
\centering
\caption{Ablation study on diagnosis candidate count $K_d$ with symptom $K_s=3$ fixed. Results show minimal variance across $K_d$ values, supporting the choice of $K_d=3$ for consistency.}
\label{tab:topk_diagnosis}
\begin{tabular}{ccc}
\toprule
\textbf{$K_d$} & \textbf{GTPA@1} & \textbf{IL} \\
\midrule
1 & 85.6\% & 19.4 \\
2 & 85.2\% & 19.5 \\
\textbf{3} & \textbf{85.2\%} & \textbf{19.4} \\
4 & 85.2\% & 19.4 \\
5 & 85.2\% & 19.4 \\
6 & 84.6\% & 19.4 \\
8 & 85.0\% & 19.4 \\
10 & 85.0\% & 19.4 \\
\bottomrule
\end{tabular}
\end{table}

% -----------------------------------------------------------------------------
% Table 4: Combined Ablation Summary
% -----------------------------------------------------------------------------
\begin{table}[t]
\centering
\caption{Summary of hyperparameter optimization. Both symptom and diagnosis selection use $K=3$ candidates based on independent ablation studies.}
\label{tab:hyperparams}
\begin{tabular}{llcl}
\toprule
\textbf{Parameter} & \textbf{Value} & \textbf{Acc.} & \textbf{Justification} \\
\midrule
Symptom $K_s$ & 3 & 85.6\% & $K \in \{1,2,3\}$ equivalent \\
Diagnosis $K_d$ & 3 & 85.2\% & 0.4\%p diff. from $K=1$ \\
Scoring & v18 & 85.6\% & Coverage-based formula \\
Min. questions & 3 & -- & Early stopping prevention \\
Max. questions & 50 & -- & Safety upper bound \\
\bottomrule
\end{tabular}
\end{table}
